\section{Example from Brenner and Foster}

\begin{frame}{Example from Brenner and Foster's Paper}
    \begin{block}{Equation 8.021 from their paper}
        \[1 + 5^a = 7^b + 7^c\]
    \end{block}
    
    \vspace{0.3cm}
    This is one of the equations they successfully solved using the modulo trick.
    
    \textbf{Why this example?}
    \begin{itemize}
        \item It's a concrete demonstration of how the method works
        \item Shows the step-by-step reasoning
        \item Illustrates the power of congruence methods
    \end{itemize}
    
    \vspace{0.2cm}
    Let's walk through their solution...
\end{frame}

\begin{frame}{Step 1: Set Up the Problem}
    We want to find all non-negative integers $(a,b,c)$ such that:
    \[1 + 5^a = 7^b + 7^c\]
    
    \vspace{0.3cm}
    \textbf{Observation:} The equation is symmetric in $b$ and $c$, so we can assume $b \leq c$ without loss of generality.
    
    \textbf{Strategy:} Use modular arithmetic to find constraints on $a$, $b$, and $c$.
\end{frame}

\begin{frame}{Step 2: Modulo 7 Analysis}
    Consider the equation modulo 7:
    \[1 + 5^a \equiv 7^b + 7^c \pmod 7\]
    
    Since $7^b \equiv 0 \pmod 7$ and $7^c \equiv 0 \pmod 7$ for $b,c \geq 1$:
    \[1 + 5^a \equiv 0 \pmod 7\]
    
    This means:
    \[5^a \equiv 6 \pmod 7\]
    
    \textbf{Question:} For which $a$ does $5^a \equiv 6 \pmod 7$?
\end{frame}

\begin{frame}{Step 3: Pattern of $5^a$ Modulo 7}
    Let's compute the pattern:
    
    \begin{tabular}{c|c}
        $a$ & $5^a \pmod 7$ \\ \hline
        1 & 5 \\
        2 & 4 \\
        3 & 6 \\
        4 & 2 \\
        5 & 3 \\
        6 & 1 \\
        7 & 5 \\
        8 & 4 \\
        9 & 6 \\
        $\vdots$ & $\vdots$
    \end{tabular}
    
    \vspace{0.2cm}
    \textbf{Pattern:} The sequence repeats every 6 steps.
    
    \textbf{Conclusion:} $5^a \equiv 6 \pmod 7$ when $a \equiv 3 \pmod 6$.
\end{frame}

\begin{frame}{Step 4: Modulo 9 Analysis}
    Now consider the equation modulo 9:
    \[1 + 5^a \equiv 7^b + 7^c \pmod 9\]
    
    First, find the pattern of $7^b$ modulo 9:
    
    \begin{tabular}{c|c}
        $b$ & $7^b \pmod 9$ \\ \hline
        1 & 7 \\
        2 & 4 \\
        3 & 1 \\
        4 & 7 \\
        5 & 4 \\
        6 & 1 \\
        $\vdots$ & $\vdots$
    \end{tabular}
    
    \vspace{0.2cm}
    So for $b \geq 1$, $7^b \in \{1,4,7\} \pmod 9$.
\end{frame}

\begin{frame}{Step 5: Compute $5^a$ Modulo 9 for $a \equiv 3 \pmod 6$}
    We know $a \equiv 3 \pmod 6$, so let's compute $5^a \pmod 9$:
    
    \begin{align*}
        5^3 &= 125 \equiv 8 \pmod 9 \\
        5^9 &= 5^{3+6} = 5^3 \cdot 5^6 \\
        5^6 &\equiv 1 \pmod 9 \text{ (by Euler's theorem)} \\
        \therefore 5^9 &\equiv 8 \cdot 1 = 8 \pmod 9
    \end{align*}
    
    In fact, for any $a \equiv 3 \pmod 6$, we have $5^a \equiv 8 \pmod 9$.
    
    Therefore:
    \[1 + 5^a \equiv 1 + 8 = 9 \equiv 0 \pmod 9\]
\end{frame}

\begin{frame}{Step 6: Compare Both Sides Modulo 9}
    We have:
    \begin{itemize}
        \item Left-hand side: $1 + 5^a \equiv 0 \pmod 9$
        \item Right-hand side: $7^b + 7^c \pmod 9$ where $7^b, 7^c \in \{1,4,7\}$
    \end{itemize}
    
    Check all possible sums:
    \begin{center}
    \begin{tabular}{c|c|c}
        $7^b$ & $7^c$ & Sum $\pmod 9$ \\ \hline
        1 & 1 & 2 \\
        1 & 4 & 5 \\
        1 & 7 & 8 \\
        4 & 4 & 8 \\
        4 & 7 & 2 \\
        7 & 7 & 5
    \end{tabular}
    \end{center}
    
    \textbf{None} of these sums equal $0 \pmod 9$!
    
    This is a contradiction unless our assumption was wrong.
\end{frame}

\begin{frame}{Step 7: Resolving the Contradiction}
    The contradiction arose because we assumed $b,c \geq 2$ when using modulo 9.
    
    Therefore, at least one of $b$ or $c$ must be $\leq 1$.
    
    \textbf{Possibilities:}
    \begin{itemize}
        \item $b = 0$ or $b = 1$
        \item $c = 0$ or $c = 1$
    \end{itemize}
    
    By symmetry, we only need to check a few cases:
    \begin{itemize}
        \item $(b,c) = (0,0)$, $(0,1)$, $(1,0)$, $(1,1)$
    \end{itemize}
    
    Only $(0,0)$ yields a solution: $(a,b,c) = (0,0,0)$.
\end{frame}

\begin{frame}{Final Result from Brenner and Foster}
    \begin{theoremss}[Brenner and Foster, 1982]
        The only non-negative integer solution to
        \[1 + 5^a = 7^b + 7^c\]
        is $(a,b,c) = (0,0,0)$.
    \end{theoremss}
    
    \vspace{0.3cm}
    \textbf{Key insight:} The modulo trick worked because:
    \begin{itemize}
        \item Modulo 7 gave a condition on $a$
        \item Modulo 9 forced a contradiction unless $b$ or $c$ was small
        \item This gave an upper bound, making the equation computable
    \end{itemize}
    
    This is a perfect example of a successful application of the modulo trick.
\end{frame}

% ================================================================================
% Section 3: The Modulo Trick - Intuitive Introduction
% ================================================================================
